\documentclass[11pt]{article}
\usepackage[T1]{fontenc}
\usepackage{lmodern}
\usepackage[margin=0.9in]{geometry}
\usepackage{amsmath,amssymb}
\usepackage{microtype}
\usepackage{booktabs}
\usepackage{xcolor}
\usepackage{enumitem}
\usepackage[colorlinks=true,linkcolor=blue!50!black,urlcolor=blue!50!black]{hyperref}

\newcommand{\C}{\mathbb{C}}
\newcommand{\abs}[1]{\lvert #1\rvert}
\newcommand{\norm}[1]{\lVert #1\rVert}
\setlist[itemize]{leftmargin=1.2em,itemsep=2pt,topsep=3pt}

\title{\bfseries Why $0.032$\\[6pt]
\large\normalfont Notes on a parameter of OpenAI, \emph{The maximum number of\\[1pt]
mutually unbiased bases in dimension six}, \S2.4 and \S3.1}
\author{}
\date{}

\begin{document}
\maketitle
\vspace{-2.6em}

\noindent
The objection is right. In my previous note the floor $N\ge47$ came from $N\ge 3/(2R_s)$, and $R_s$
was assumed. That is a relation between two free parameters, not a derivation of either. The preprint
does not derive them either: it says, in as many words, ``Set $R=0.032$'' (and, earlier, ``Set
$N=112$''). What can be done is to say what the two numbers \emph{are}, which quantities each one
bounds, and why each is squeezed from both sides. Then the honest conclusion is that $(N,R_s)$ is a
joint budget whose feasible set is a rectangle, and $(112,0.032)$ is one point in it.

\section*{1.\ What $0.032$ measures}

The seed ball is a statement of \emph{uncertainty}: Proposition 2.10 certifies that each of the nine
seed phases of an exact Hadamard matrix lies within $0.032$ cycles of some recorded centre, that is,
within
\[
R_s=2\pi r_s=2\pi(0.0320001)=0.201
\]
radians. Section 3 then spends that uncertainty, once, in the chord bound: writing $E',B',C'$ for the
seed blocks of an exact matrix and $E_0,B_0,C_0$ for the centre's,
\begin{equation}\label{eq:chord}
e_1^2+(b_1')^2+(c_1')^2=\norm{\Delta E}_F^2+\norm{\Delta B}_F^2+\norm{\Delta C}_F^2\le\norm{s}^2\le R_s^2 ,
\end{equation}
by the chord bound $\abs{e^{ix}-1}\le\abs{x}$. Every later estimate in Stages A and B is a value
computed at the centre \emph{plus a multiple of $R_s$}; the cleanest instance is
\begin{equation}\label{eq:Astar}
A_0=\norm{E_0}_{\mathrm{op}}=\sqrt{2+\abs{1+a}}\,,
\qquad A_*=\min\bigl(2,\;A_0+R_s\bigr),
\end{equation}
where $A_*$ bounds $\norm{E'}_{\mathrm{op}}$ for the exact matrix: the displacement costs $R_s$, and
the alternative bound $2$ holds for any $2\times2$ matrix of unit entries.

So $R_s$ is not a resolution parameter like $N$; it is an \emph{error budget in radians}, and the
question ``why $0.032$'' is the question ``why is $0.2$ radians of uncertainty affordable downstairs,
and why is less of it unaffordable upstairs''.

\section*{2.\ What bounds $R_s$ from above: the checks must stay conclusive}

Each check in Stage A has a threshold, and each threshold is approached by the exact matrix at a rate
proportional to $R_s$. Three examples, all from \S3:
\begin{itemize}
\item \emph{Frame conditioning.} The frame built from a unit $4$-vector $p$ has squared singular
values $4\pm\abs{S}$, $S=\sum_jp_j$, and the pseudo-inverse that compresses the unknown block is
\[
M(M^*M)^{-1}=\frac{1}{16-\abs{S}^2}\bigl(4\mathbf 1-p\overline S,\;4p-\mathbf 1S\bigr).
\]
and the check needs $\abs{S}<4$ with a margin: the denominator is small exactly when the frame is
nearly degenerate, and a displacement of $R_s$ moves $S$ by $O(R_s)$.
\item \emph{Rank.} With $m_B=\sigma_{\min}(B_0)$, $m_C=\sigma_{\min}(C_0)$ and $m=\min(m_B,m_C)>0$,
the projections $P_B=I-G_BB_0$ need $m>0$ to be defined at all.
\item \emph{Compression.} The trace defect $q$ of $6I-Y^*Y$ and the distance $p\le q/(\sqrt6+\sqrt{6-q})$
of the compressed block from $\sqrt6\,U(2)$ both grow with $R_s$; the paper's parameter range has
$q<0.32$ and $p\lesssim0.09$.
\end{itemize}
None of these yields a closed-form ceiling on $R_s$. Each says: \emph{for every retained centre, $R_s$
must be smaller than that centre's margin divided by the constant with which the displacement enters
it}. The admissible maximum is therefore the minimum over some hundreds of thousands of centres, which
is a computed number, not a derived one. It is also the reason the paper checks each centre
individually and keeps a failed centre \emph{unresolved} instead of excluding it: a failed check is
evidence about $R_s$, never a proof about Hadamards.

\section*{3.\ What bounds $R_s$ from below: the cover must be affordable}

The seed set is not nine-dimensional: a dephased Hadamard matrix has $25$ phases and the orthogonality
conditions are $21$ real equations, so off the singular loci the set is \emph{four}-dimensional.
Covering a four-dimensional set by balls of radius $R_s$ costs about $R_s^{-4}$ balls:

\begin{center}
\begin{tabular}{@{}ccccc@{}}
\toprule
$R_s$ (cycles) & seeds $\approx R_s^{-4}$ & smallest $N$ with $\tfrac32w\le R_s$ & boxes $\binom{N+4}{5}$ at that $N$\\
\midrule
$0.064$ & $6.0\times10^{4}$ & $24$ & $9.8\times10^{4}$\\
$\mathbf{0.032}$ & $\mathbf{9.5\times10^{5}}$ & $\mathbf{47}$ & $\mathbf{2.3\times10^{6}}$\\
$0.016$ & $1.5\times10^{7}$ & $94$ & $6.8\times10^{7}$\\
\bottomrule
\end{tabular}
\end{center}

\noindent
Halving $R_s$ makes the seed cover sixteen times dearer and forces bins twice as fine. This keeps $R_s$
from being tiny: the seed radius is what the cover can afford, the bin width is what the enumeration
can afford, and $\tfrac32w\le R_s$ is the only inequality that couples them.

\section*{4.\ So what is $0.032$?}

It is the point where two costs of opposite sign meet, and the meeting point is found by measurement.
\emph{Downstairs}, a smaller $R_s$ costs balls like $R_s^{-4}$ and forces finer bins like $R_s^{-1}$,
so the first stage becomes unaffordable. \emph{Upstairs}, a larger $R_s$ inflates every Stage A
bound by a multiple of it --- frame denominators, rank margins, $q$, $p$ --- and every Stage B and C
ball downstream is measured against $R_s$ as well, so the per-seed work grows with $R_s$. The paper
reports a feasibility check (Section 7) that fails if any centre remains unresolved; $0.032$ is a
value at which it succeeds, and the ``why'' is that no larger value kept the checks conclusive and no
smaller one kept the cover affordable.

Two consequences worth stating. First, $0.032$ is not a threshold the mathematics depends on: any
radius in the feasible window works, with Section A's constants recomputed. Second --- the correction
to my earlier note --- $(112,\,0.032)$ is neither minimal nor canonical. At $R_s=0.032$ the coupling
permits $N=47$, which would cut the enumeration from $1.60\times10^{8}$ boxes to $2.3\times10^{6}$,
seventy times cheaper. The paper spends that factor on bins $2.4$ times finer than the coupling
requires, because finer seeds mean smaller Taylor remainders, fewer $U(2)$ cells per seed and fewer
vector balls in Stage C. That is tuning, not theorem, and one can disagree with it without touching the
proof.

\medskip
\noindent
The lesson is worth more than either number. There are two independent budgets --- how finely the
space is cut, and how much error each cut is allowed to carry --- coupled by one geometric inequality
(a box must fit in a ball). Neither follows from the other, and neither follows from the conjecture;
both are chosen inside a window and then \emph{frozen}, because Section A is proved for those values.
Asking ``why 112'' or ``why 0.032'' is asking for the window, and the window is the honest answer.

\end{document}
